\section{Introduction}
\label{sec:intro}




\begin{figure}[t]
    \centering
    \includegraphics[width=0.90\linewidth]{Figures/L_system.pdf}
    \caption{\revised{Intelligence levels of video analytics systems.}}
    \label{fig:l_system}
\end{figure}

Video analytics~\cite{remix_mobicom_21, bhardwaj2022ekya, khani2023recl, 10.1145/3624478} has emerged as a transformative technology across a wide array of domains, such as environment monitoring, intelligent transportation systems, industrial automation, and retail monitoring. \revised{By leveraging the capabilities of deep learning (DL) models, video analytics systems are able to extract patterns, derive meaningful insights, and generate actionable information from video data, thereby facilitating more efficient and precise monitoring, detection, and response to events. The desired features of video analytics systems necessitate a variety of capabilities. As shown in Fig.~\ref{fig:l_system}, We categorize the intelligence levels of both current and prospective video analytics systems into five tiers, designated as L1 through L5:}

\textbf{L1} systems for specific classification, segmentation, and detection using models \eg ResNet~\cite{he2016resnet} and EfficientDet~\cite{tan2020efficientdet} to extract spatial information from video data, including object classes and bounding boxes, etc~\cite{remix_mobicom_21, bhardwaj2022ekya, khani2023recl}.


\textbf{L2} systems go beyond spatial information extraction by enabling causal event detection and analytics, \ie identifying short-term events. They use models like C3D~\cite{tran2015c3d} and ActionFormer~\cite{zhang2022actionformer} to detect and localize events (\eg actions, activities, anomalies) through spatiotemporal modeling.




\revised{\textbf{L3} systems advance beyond L2’s spatiotemporal detection by integrating neural language processing (NLP) capabilities. Leveraging models such as CLIPBERT~\cite{lei2021less}, these systems are able to interpret and respond to natural language queries \eg \textit{"What animals appeared in the videos?"} rather than simply detecting (manually) predefined events. Although query handling remains confined to specific domains, L3 systems significantly enhance user interaction and accessibility in video analytics.}






Despite significant advancements, current video analytics systems~\cite{khani2023recl, bhardwaj2022ekya, zhang2024vulcan, padmanabhan2023gemel, lu2022turbo, sivaraman2024gemino, zhang2021sensei, remix_mobicom_21} primarily focus on L1 to L3 levels. These systems are designed for \emph{closed-end} analytics, relying on specialized models for specific tasks, which limits their flexibility and adaptability. Consequently, we envision L4 and L5 systems to enable \emph{open-end} analytics.



\revised{\textbf{L4} systems represent a significant leap forward by enabling open-ended video comprehension, reasoning, and analytics. These systems support general-purpose video analytics, effectively processing natural language queries and generating contextually appropriate responses, while also facilitating complex, long-term spatiotemporal reasoning. For instance, L4 systems can address inquiries such as: \textit{"What abnormal events occurred in the past ten hours?"}, \textit{"What caused the person to fall?"}, or \textit{"How did the animals behave after appearing on camera?"}}


\revised{\textbf{L5} systems advance beyond L4 by autonomously integrating external public and domain-specific knowledge sources to uncover both explicit and implicit relationships between video content and broader world knowledge. This capability fulfills the ultimate objective of video analytics: deriving profound insights and delivering actionable, automated solutions.}


In this paper, we delve into the development of \textbf{L4 video analytics systems}, leveraging the transformative potential of vision language models (VLMs). By combining vision and language understanding, VLMs enable generalized visual detection and advanced video comprehension, including causal reasoning, key information retrieval, and human-interpretable explanations. The integration of perception, reasoning, and interaction makes VLMs highly adaptable, positioning them as a key foundation for L4 systems.


However, integrating VLMs into video analytics poses significant challenges, primarily due to the limited context window of current VLMs compared to the extensive duration of video sources in typical video analytics scenarios. While L1 to L3 systems handle spatial and short-term causal event detection by processing frames independently or within small sliding windows (\eg a few seconds), L4 systems require collective analysis of related frames for long-term causal detection, summarization, and reasoning. Current VLMs, like QwenVL~\cite{wang2024qwen2}, can process up to 768 frames, covering minutes or hours of video. However, video analytics often involves much longer sources, spanning  hundreds of hours or continuous streams, far beyond the capabilities of existing VLMs.

Recent studies~\cite{shen2024longvu, shen2025longvita, wang2025adaretake} that attempt to extend the context window of VLMs remain inadequate for processing video sources spanning hundreds of hours. Retrieval-augmented generation (RAG) frameworks~\cite{edge2024graphrag, guo2024lightrag, fan2025minirag, lazygraphrag2024} aim to address similar limitations by first retrieving relevant frames from massive contents and then generating final answers. However, these approaches still face significant challenges in handling the video modality (as detailed in \S\ref{sec:evaluation}), leading to both reduced analytics accuracy and substantial computational overhead.


In this paper, we present \sysname, a system that integrates VLMs into video analytics to enable L4 capabilities. The core innovation of \sysname is its efficient indexing mechanism, designed to handle extremely long video sources or unlimited video streams, and by leveraging the index, \sysname effectively retrieves relevant information for a given query and generates accurate, robust responses. Specifically, \sysname introduces the following key features that distinguish it from existing systems: 1) analyzing extremely long videos, spanning hundreds of hours or even unlimited video streams; 2) supporting near-real-time (\eg at more than 1 FPS) index construction and analytics; and 3) handling diverse and complex queries including temporal grounding and reasoning, summarization, event and entity understanding, and key information retrieval.




Particularly, \sysname introduces two key components: \emph{near-real-time index construction} and \emph{agentic retrieval and generation}. During the index construction phase (\S\ref{sec:offline}), we propose event knowledge graphs (EKGs) as an indexing mechanism for video analytics. Unlike traditional knowledge graphs (KGs) used in text-based RAG systems~\cite{edge2024graphrag, guo2024lightrag, lazygraphrag2024}, EKGs represent a flow of insightful events, effectively capturing video dynamics and temporal consistency. \revised{Using a small VLM, such as Qwen2.5-VL-7B}, \sysname extracts information to construct the EKG. To optimize this process, we introduce techniques that enable near-real-time index construction, achieving more than 5 FPS on typical edge servers equipped with 2 $\times$ 4090 GPUs.


In the retrieval and generation phase (\S\ref{sec:online}), \sysname employs an agentic search mechanism instead of directly retrieving information from the constructed index. This approach allows \sysname to proactively retrieve more relevant information by utilizing contextual hints captured within the EKG, enables \sysname to handle complex queries, including summarization and multi-hop reasoning etc. Furthermore, we introduce techniques to enhance \sysname to robustly generate answers based on the retrieved information.



We evaluate \sysname on multiple public video understanding benchmarks, including LVBench~\cite{wang2024lvbench} and VideoMME-Long~\cite{fu2024videomme}. These benchmarks collectively comprise approximately 400 videos and 2,500 questions. We compare \sysname against a diverse range of baselines, \revised{including state-of-the-art (SOTA) VLMs such as GPT-4o~\cite{achiam2023gpt}, Gemini-1.5-pro}~\cite{team2023gemini}, Phi-4-Multimodal~\cite{abdin2024phi3}, Qwen2.5-VL-7B~\cite{wang2024qwen2}, InternVL2.5-8B~\cite{chen2024internvl}, and LLaVA-Video-7B~\cite{zhang2024llavavideo}, as well as typical video-RAG systems such as VideoTree~\cite{wang2024videotree}, VideoAgent~\cite{wang2024videoagent}, DrVideo~\cite{ma2024drvideo}, and VCA~\cite{yang2024vca}. On both benchmarks, \sysname establishes new SOTA performance, achieving 62.3\% on LVBench and 64.1\% on VideoMME-Long, respectively, significantly surpassing baselines by up to 16.9\% and 5.2\%.




In addition to the public video understanding benchmarks, we introduce a new benchmark, \sysname-100, specifically designed to evaluate L4 video analytics tasks. \sysname-100 comprises 8 ultra-long videos, each exceeding 10 hours in duration, and includes a total of 120 manually annotated questions and answers. The videos are carefully curated from typical video analytics scenarios, and the questions cover multiple key categories relevant to video analytics systems. Evaluation results show that \sysname achieves significantly better performance on \sysname-100 compared to various baselines, with improvements of approximately 20.8\%. In summary, we make the following contributions in this paper:




\begin{itemize}
    \item We propose \sysname, the first L4 video analytics system powered by VLMs, to the best of our knowledge.
    \item \sysname introduces near-real-time index construction and agentic retrieval and generation, along with innovative techniques that enable key features for L4 video analytics, including open-ended analytics on extremely long video sources in near-real-time.  
    \item We evaluate \sysname on two widely-used video understanding benchmarks, achieving SOTA performance with 62.3\% on LVBench and 64.1\% on VideoMME-Long.  
    \item Furthermore, we present \sysname-100, a benchmark specifically designed for L4 video analytics systems, where \sysname demonstrates significant improvements, outperforming baselines by approximately 20.8\%.  
\end{itemize}


















